% =============================================================================
% PROTEA — Documento de proyecto
% Versión para revisión interna (David Orellana)
% Compila con pdflatex: pdflatex PROTEA_project_document.tex
% =============================================================================

\documentclass[11pt,a4paper]{article}

\usepackage[utf8]{inputenc}
\usepackage[T1]{fontenc}
\usepackage[spanish]{babel}
\usepackage[margin=2.5cm]{geometry}
\usepackage{microtype}
\usepackage{parskip}
\usepackage{enumitem}
\usepackage{xcolor}
\usepackage{hyperref}
\usepackage{booktabs}
\usepackage{array}
\usepackage{graphicx}
\usepackage{titlesec}

\hypersetup{
  colorlinks=true,
  linkcolor=black,
  urlcolor=blue!60!black,
  citecolor=black,
  pdftitle={PROTEA — Documento de proyecto},
  pdfauthor={Francisco Miguel Pérez Canales}
}

\titleformat{\section}{\Large\bfseries}{\thesection}{1em}{}
\titleformat{\subsection}{\large\bfseries}{\thesubsection}{1em}{}

\title{\textbf{PROTEA}\\[0.3em]
  \large Plataforma abierta de anotación funcional de proteínas\\
  basada en modelos de lenguaje y re-ranking explicable\\[0.8em]
  \normalsize Documento de proyecto — versión para revisión interna}

\author{Francisco Miguel Pérez Canales\\
  \small Universidad de Sevilla\\
  \small \texttt{frapercan1@alum.us.es}}

\date{Abril 2026}

\begin{document}

\maketitle

\begin{abstract}
\noindent
PROTEA es una plataforma abierta para la predicción de función proteica mediante
modelos de lenguaje de proteínas, búsqueda por vecinos más cercanos (KNN) sobre
embeddings vectoriales, y un re-ranker explicable que incorpora información de
alineamiento y taxonomía. El proyecto responde a dos objetivos que la comunidad
de anotación funcional todavía no ha resuelto conjuntamente: (1) obtener
predicciones competitivas con el estado del arte en el benchmark de referencia
(CAFA), y (2) hacerlo desde una plataforma reproducible, trazable y usable por
investigadores sin perfil técnico. PROTEA se sitúa actualmente en el segundo
puesto del \emph{leaderboard} público de CAFA~6 (a la espera del cierre oficial
del \emph{ranking} privado) y está preparado para ser desplegado como servicio
estable. Este documento describe el enfoque metodológico, la arquitectura
técnica, el estado actual del proyecto, los recursos necesarios para el
despliegue y la hoja de ruta de publicación. Se entrega como base para
discusión y revisión interna con colaboradores.
\end{abstract}

\tableofcontents

\vspace{1em}

% =============================================================================
\section{Motivación}
% =============================================================================

La anotación funcional de proteínas —la asignación de términos de Gene Ontology
(GO) a secuencias proteicas— es un problema central en biología computacional.
A medida que la secuenciación masiva produce proteomas de organismos no-modelo
a un ritmo que supera con mucho la capacidad de anotación experimental, la
predicción automática de función se ha convertido en un cuello de botella real
para la biología de biodiversidad, la metagenómica y la genómica comparada.

Las herramientas clásicas del campo (InterProScan, eggNOG-mapper, PANNZER2)
dependen fundamentalmente de búsquedas basadas en similitud de secuencia,
modelos ocultos de Markov (HMMs) y ortología. Estos métodos funcionan bien en
organismos bien representados en las bases de datos de referencia, pero pierden
potencia cuando se aplican a especies no-modelo, proteomas recién ensamblados,
o dominios de secuencia divergentes.

Desde 2020, los modelos de lenguaje de proteínas (ESM, ProstT5, Ankh, etc.) han
demostrado que los embeddings vectoriales derivados de estos modelos capturan
información estructural y funcional que no es accesible a los métodos clásicos.
Sin embargo, la traducción de esa capacidad en herramientas prácticas para la
comunidad sigue siendo desigual: la mayoría de sistemas publicados son
\emph{scripts} de investigación que requieren conocimiento técnico, no ofrecen
interfaces usables y no garantizan reproducibilidad bit a bit entre
ejecuciones.

PROTEA se diseña para cerrar esa brecha: predicciones competitivas con el
estado del arte, expuestas a través de una plataforma web reproducible,
auditable y usable sin conocimientos de línea de comandos.

% =============================================================================
\section{Estado del arte}
% =============================================================================

El estándar de evaluación del campo es la competición internacional CAFA
(\emph{Critical Assessment of protein Function Annotation}), que se celebra
aproximadamente cada tres años y que establece el \emph{benchmark} público más
riguroso para predicción de función. Se organiza sobre un \emph{holdout}
temporal de anotaciones experimentales: los sistemas predicen función para un
conjunto de proteínas cuya anotación real todavía no se conoce públicamente en
el momento del envío, y se evalúan meses después cuando las anotaciones se
hacen públicas.

Los participantes de referencia en CAFA 5 y CAFA 6 incluyen:

\begin{itemize}[leftmargin=1.5em]
  \item \textbf{Herramientas clásicas} basadas en búsqueda de similitud y
    anotación por transferencia: InterProScan (consorcio InterPro, basado en
    perfiles HMM de múltiples bases), eggNOG-mapper (anotación por ortología a
    partir de la base eggNOG) y PANNZER2 (predicción bayesiana sobre alineamientos
    BLAST+ con ponderación de evidencias).
  \item \textbf{Sistemas basados en aprendizaje profundo sobre secuencias}:
    DeepGOPlus, NetGO~2.0/3.0, GOPredSim, entre otros. Estos sistemas combinan
    embeddings, redes convolucionales o modelos basados en atención con
    clasificadores supervisados sobre el vocabulario GO.
  \item \textbf{Sistemas mixtos} que combinan evidencias heurísticas de
    alineamiento, ortología y embeddings, típicamente ensamblados mediante
    \emph{stacking} manual.
\end{itemize}

El trabajo presentado en este documento se posiciona en una tercera vía: usar
embeddings de proteínas como representación primaria, pero sustituir el
clasificador supervisado por una búsqueda KNN sobre un corpus de referencia
anotado, y a continuación aplicar un re-ranker explicable que re-ordena los
candidatos incorporando información de alineamiento y taxonomía. Esta
combinación preserva la capacidad discriminativa de los modelos de lenguaje y
añade una capa de interpretabilidad de la que los clasificadores monolíticos
carecen.

% =============================================================================
\section{Enfoque metodológico}
% =============================================================================

\subsection{Representación: embeddings reproducibles}

Cada secuencia del corpus de referencia se codifica en un vector mediante un
modelo de lenguaje de proteínas. PROTEA soporta actualmente cuatro backends
(ESM-C, ProstT5, Ankh y T5 clásico) con normalización unificada de tokens
especiales, de forma que los índices posicionales en los embeddings coinciden
exactamente con las posiciones de aminoácido.

La novedad práctica a nivel de sistema es que cada \emph{receta} de embedding
—modelo, versión, \emph{chunking}, precisión numérica— se materializa como un
objeto \texttt{EmbeddingConfig} identificado por UUID, que actúa como clave
primaria de los vectores almacenados en la base de datos. Esto garantiza
reproducibilidad bit a bit: regenerar dos veces los embeddings de la misma
secuencia bajo el mismo \texttt{EmbeddingConfig} produce vectores idénticos.

\subsection{Búsqueda KNN sobre corpus anotado}

Dada una secuencia de consulta embebida, PROTEA recupera los $k$ vecinos más
cercanos en el corpus de referencia mediante búsqueda por similitud vectorial.
El corpus de referencia se construye a partir de instantáneas de Gene Ontology
Annotation (GOA) y QuickGO, con filtrado por fecha para permitir evaluaciones
con \emph{holdout} temporal estricto (críticas para el realismo del
\emph{benchmark}).

Cada vecino contribuye su conjunto de términos GO asignados al ranking de
predicciones de la consulta, ponderado inicialmente por distancia en el espacio
de embeddings.

\subsection{Re-ranker con información estructural y taxonómica}

La mejora clave respecto a un KNN puro consiste en aplicar un re-ranker
posterior que incorpora dos fuentes de información adicionales:

\begin{enumerate}[leftmargin=1.5em]
  \item \textbf{Alineamiento de secuencia}: para cada par (consulta, vecino) se
    calcula un alineamiento global (Needleman-Wunsch) y local (Smith-Waterman)
    mediante la librería \texttt{parasail} con matriz BLOSUM62, extrayendo
    identidad, similitud, cobertura, huecos y longitud del alineamiento.
  \item \textbf{Información taxonómica}: se extrae la distancia taxonómica
    entre la especie de la consulta y la de la referencia (LCA, distancia,
    ancestros comunes, relación taxonómica clasificada) usando la taxonomía
    NCBI.
\end{enumerate}

El re-ranker es un modelo supervisado (\emph{gradient boosting}) que combina
todas las \emph{features} anteriores —distancia en embeddings, métricas de
alineamiento, información taxonómica— para producir un \emph{score} final por
término GO candidato. Este planteamiento tiene dos ventajas claras sobre los
sistemas basados exclusivamente en embeddings:

\begin{itemize}[leftmargin=1.5em]
  \item \textbf{Explicabilidad intrínseca}: cada predicción puede justificarse
    mostrando al usuario la proteína de referencia que la motivó, el
    alineamiento que comparte con la consulta, y la relación taxonómica entre
    ambas. Esto es especialmente relevante en anotación funcional, donde el
    biólogo final necesita poder evaluar la credibilidad de la predicción
    antes de aceptarla.
  \item \textbf{Robustez ante \emph{drift} de dominio}: cuando la consulta
    pertenece a un organismo mal representado, el re-ranker puede compensar
    mediante las \emph{features} de alineamiento, recuperando capacidad
    predictiva donde un clasificador monolítico falla silenciosamente.
\end{itemize}

\subsection{Entrenamiento del re-ranker sobre \emph{holdout} temporal}

El re-ranker se entrena sobre parejas generadas a partir de instantáneas
históricas de GOA: se usa una ventana de lanzamientos anteriores como
\emph{train} y un lanzamiento posterior como \emph{holdout} (actualmente
GOA~160–225 para entrenamiento, GOA~230 como holdout). Este protocolo imita
exactamente el escenario de evaluación de CAFA —predecir anotaciones que
todavía no se conocen en el momento de entrenamiento— y garantiza que los
resultados medidos en CAFA no están inflados por fugas temporales.

% =============================================================================
\section{Arquitectura técnica}
% =============================================================================

PROTEA se implementa como una plataforma completa que incluye base de datos,
sistema de colas, API REST e interfaz web. La arquitectura está organizada en
torno a tres pilares de diseño que merecen mención explícita porque son los
que diferencian PROTEA de las herramientas clásicas del campo y los que van a
estructurar la eventual publicación de ingeniería.

\subsection{Pilar 1 — Embeddings versionados como ciudadanos de primera clase}

En la mayoría de herramientas bioinformáticas que usan embeddings, los
vectores se almacenan en archivos sueltos, con nombres basados en convenciones
informales, y la \emph{receta} usada para generarlos queda implícita en el
código que los produjo. Cuando se reproduce un experimento meses después, a
menudo es imposible saber con certeza con qué modelo, con qué parámetros y con
qué versión de la librería se generaron los vectores originales.

PROTEA promueve la configuración de embedding a entidad de primera clase en la
base de datos (\texttt{EmbeddingConfig}), identificada por UUID y referenciada
como clave foránea por cada vector almacenado. Esto permite, entre otras
cosas:

\begin{itemize}[leftmargin=1.5em]
  \item Regenerar deterministamente cualquier vector a partir de su
    \texttt{EmbeddingConfig} y la secuencia original.
  \item Comparar predicciones entre configuraciones distintas sin ambigüedad
    sobre qué se está comparando exactamente.
  \item Migrar entre precisiones numéricas (por ejemplo, de \texttt{fp32} a
    \texttt{halfvec}/\texttt{fp16}) con trazabilidad completa del cambio.
\end{itemize}

\subsection{Pilar 2 — Operaciones como funciones puras del dominio}

La lógica de dominio de PROTEA se organiza en torno al protocolo
\texttt{Operation}, que exige que cada unidad de trabajo implemente una firma
simple: un nombre y una función \texttt{execute(session, payload, *, emit)}
que recibe una sesión de base de datos, la carga útil del trabajo y un
\emph{callback} para reportar progreso. Las operaciones no gestionan conexiones
a RabbitMQ, no abren sesiones de base de datos por su cuenta, y no conocen la
infraestructura de colas subyacente.

Esta separación es una aplicación directa de los principios de \emph{hexagonal
architecture}~\cite{cockburn2005hexagonal} y \emph{clean
architecture}~\cite{martin2017clean} al dominio del \emph{research software}
bioinformático, donde la norma histórica han sido \emph{workers} monolíticos
que acoplan SQL, AMQP, lógica de dominio e inferencia en clases de varios
cientos de líneas. No se reclama aquí la invención del patrón —existe en
ingeniería de software desde al menos 2005— sino su aplicación sistemática a
un campo que tradicionalmente no lo usa, y la demostración empírica de que
esa disciplina habilita testabilidad sin infraestructura, extensibilidad,
reutilización y reproducibilidad que los sistemas monolíticos no ofrecen.

\subsection{Pilar 3 — Sistema de colas heterogéneo con coordinación de dominio}

El tercer pilar es donde se concentra la contribución técnica original de
PROTEA. El sistema de colas enruta trabajos a siete colas distintas según su
naturaleza, y combina dos estilos de \emph{worker} en el mismo \emph{framework}:

\begin{itemize}[leftmargin=1.5em]
  \item \textbf{\texttt{QueueConsumer}}, que gestiona transiciones de estado y
    filas de \texttt{Job} completas en la base de datos, apropiado para
    trabajos de larga duración que requieren auditoría.
  \item \textbf{\texttt{OperationConsumer}}, que procesa cargas útiles directas
    sin crear filas de \texttt{Job}, apropiado para micro-trabajos de alto
    volumen (por ejemplo, cada \emph{batch} de inferencia dentro de un
    \emph{job} mayor).
\end{itemize}

Sobre estos dos consumidores se implementan tres patrones que, por lo que
conocemos, no están presentes en la literatura de \emph{research software}
bioinformático:

\begin{enumerate}[leftmargin=1.5em]
  \item \textbf{Coordinadores con \emph{backpressure} de dominio.} Las
    operaciones que dependen de recursos escasos (típicamente la GPU) levantan
    una excepción \texttt{RetryLaterError} cuando el recurso no está
    disponible; el coordinador atrapa la excepción, re-encola el trabajo con
    retardo y libera el \emph{worker}. Esto expresa la contrapresión a nivel
    de dominio, no de infraestructura, y permite serialización natural del
    acceso a GPU sin necesidad de semáforos externos ni \emph{locks}
    distribuidos.
  \item \textbf{Jerarquía \emph{parent-child} con progreso atómico.} Un trabajo
    padre (por ejemplo, \emph{calcular embeddings para un proteoma de 50{.}000
    secuencias}) se descompone en centenares de trabajos hijos que no crean
    fila de \texttt{Job} propia, sino que incrementan atómicamente el contador
    de progreso del padre. Esto reduce el volumen de escritura en la tabla de
    \texttt{Job} en uno o dos órdenes de magnitud sin perder visibilidad del
    progreso.
  \item \textbf{Despliegue federado de \emph{workers} bajo demanda.} Los
    \emph{workers} no necesitan vivir en la misma máquina que el plano de
    control (base de datos, cola, API): basta con que tengan acceso de red al
    RabbitMQ central. Esto habilita topologías híbridas donde el plano de
    control reside en un servicio persistente (VM, \emph{cloud}, servidor
    institucional) y los \emph{workers} pesados se lanzan como \emph{jobs} en
    infraestructuras externas (HPC con SLURM, nodos \emph{cloud} efímeros,
    estaciones de trabajo con GPU) sin tocar una línea del plano de control.
\end{enumerate}

Este tercer pilar es, en términos de contribución original, el más fuerte del
proyecto, y el que justifica la aspiración a publicar la parte de ingeniería
en una revista de alto impacto.

% =============================================================================
\section{Estado actual del proyecto}
% =============================================================================

\subsection{Resultados}

\begin{itemize}[leftmargin=1.5em]
  \item \textbf{Posición en CAFA~6}: actualmente \#2 en el \emph{leaderboard}
    público, a la espera del cierre oficial del \emph{ranking} privado.
  \item \textbf{Comparación con el estado del arte}: se han obtenido métricas
    bajo el mismo marco de evaluación (el oficial de CAFA) para InterProScan,
    eggNOG-mapper y PANNZER2 sobre el mismo conjunto de validación, lo que
    permite una comparación directa y justa. Los resultados preliminares
    sostienen que el enfoque basado en embeddings + re-ranker supera
    claramente a los métodos basados en HMM y ortología en los tres aspectos
    de GO (BPO, MFO, CCO). A diferencia de las aproximaciones habituales, el
    mecanismo de evaluación y validación está \emph{integrado de forma
    transparente en la propia plataforma}: el marco CAFA se ejecuta como parte
    del \emph{pipeline} sobre las mismas estructuras de datos y el mismo flujo
    de trabajo que la inferencia, sin \emph{scripts} externos ni post-proceso
    manual, garantizando que cada conjunto de predicciones es reproducible y
    comparable extremo a extremo.
  \item \textbf{Plataforma desplegable}: código base estable, con una
    \emph{suite} de tests amplia y cobertura notablemente por encima del
    umbral habitual en \emph{research software}, documentación Sphinx al día y
    \emph{pipeline} completo funcionando desde la carga de un FASTA por el
    usuario hasta la descarga del TSV de predicciones.
\end{itemize}

\subsection{Infraestructura de software}

\begin{itemize}[leftmargin=1.5em]
  \item \textbf{Stack}: FastAPI (API REST), SQLAlchemy 2.0 (ORM), PostgreSQL 16
    con extensión pgvector, RabbitMQ (sistema de colas), Next.js 16 con
    Tailwind v4 (interfaz web).
  \item \textbf{Orquestación}: \texttt{docker compose} para desarrollo y
    despliegue local; imágenes publicadas automáticamente en \emph{ghcr.io} en
    cada \emph{push} a \texttt{main}.
  \item \textbf{Gestión de procesos}: \emph{script} \texttt{manage.sh} que
    arranca, para, escala y monitoriza los 9 procesos (API, \emph{workers}
    especializados, \emph{frontend}) en desarrollo local.
\end{itemize}

\subsection{Trabajo pendiente a corto plazo}

\begin{itemize}[leftmargin=1.5em]
  \item Iteración final limpia de la experimentación con traza completa de
    hiperparámetros, \emph{datasets} y reproducibilidad verificable.
  \item Reconstrucción del método que obtuvo la posición en CAFA~6 en un
    \emph{pipeline} documentado y reproducible de principio a fin (requisito
    imprescindible para la publicación del método).
  \item Ajustes finales de documentación y \emph{tutorials}.
  \item Publicación de \emph{preprint} en bioRxiv como paso previo al envío a
    revista.
\end{itemize}

% =============================================================================
\section{Infraestructura de despliegue}
% =============================================================================

\subsection{Situación actual: despliegue doméstico}

En el momento de redactar este documento, PROTEA se ejecuta íntegramente
sobre el equipo personal del autor, expuesto públicamente mediante un túnel
\emph{ngrok} sobre un subdominio reservado. Toda la plataforma —API,
PostgreSQL con \emph{pgvector}, RabbitMQ, \emph{workers} de cómputo,
interfaz web y documentación Sphinx— corre en esta única máquina. El
\emph{stack} se arranca con \texttt{bash scripts/manage.sh start} y el túnel
con un \emph{script} auxiliar, quedando accesible desde Internet sin
intervención de un administrador de sistemas.

Las características del equipo son:

\begin{center}
\begin{tabular}{ll}
\toprule
\textbf{Componente} & \textbf{Descripción} \\
\midrule
CPU     & Intel Core i5-12400F (6 núcleos / 12 \emph{threads}) \\
RAM     & 64 GB DDR4 \\
GPU     & NVIDIA GeForce RTX 3060, 12 GB VRAM \\
Disco   & NVMe 1 TB \\
Sistema & Ubuntu 24.04 LTS (kernel Linux 6.17) \\
Red     & Fibra doméstica + túnel \emph{ngrok} con subdominio reservado \\
\bottomrule
\end{tabular}
\end{center}

El uso observado durante las cargas de trabajo habituales —generación de
\emph{embeddings}, inferencia, \emph{benchmarking} del re-ranker— es
intensivo pero sostenible: la GPU se mantiene por debajo del 90\,\% de
utilización y dentro del rango térmico seguro (del orden de
65\,\textdegree{}C durante la inferencia, lejos del umbral de \emph{throttling}),
la CPU estabiliza el \emph{package} en torno a 60\,\textdegree{}C bajo carga
continua, el almacenamiento NVMe se mantiene a $\sim$33\,\textdegree{}C y la
ocupación del disco ronda el 70\,\%. En otras palabras, la máquina es
capaz de ejecutar la plataforma extremo a extremo, pero opera cerca del
techo razonable para un equipo doméstico no diseñado para servicio 24/7:
cualquier uso externo sostenido requiere trasladar el servicio a una
máquina de alojamiento estable.

\subsection{Requisitos}

Para desplegar PROTEA como servicio accesible se necesita una máquina
—física o virtual— con las siguientes características mínimas:

\begin{center}
\begin{tabular}{lll}
\toprule
\textbf{Recurso} & \textbf{Mínimo} & \textbf{Recomendado} \\
\midrule
vCPU        & 4           & 8 \\
RAM         & 16 GB       & 32 GB \\
Disco SSD   & 500 GB      & 1 TB \\
Sistema     & \multicolumn{2}{l}{Linux (Ubuntu 22.04 / Debian 12), Docker, \texttt{docker compose}} \\
Red         & \multicolumn{2}{l}{IP pública o subdominio HTTPS, puertos 80/443 accesibles} \\
Acceso      & \multicolumn{2}{l}{SSH administrativo, salida a Internet (descargas UniProt / GO)} \\
GPU         & \multicolumn{2}{l}{\textbf{No requerida} (ver más abajo)} \\
\bottomrule
\end{tabular}
\end{center}

\subsection{Despliegue federado y ausencia de requisito GPU}

La arquitectura de PROTEA separa deliberadamente el \emph{plano de control}
—base de datos, sistema de colas, API, interfaz web, todos persistentes y de
consumo moderado— del \emph{plano de datos} —\emph{workers} de inferencia con
GPU, efímeros y desplegables bajo demanda en infraestructuras externas. La
máquina de alojamiento no necesita GPU.

Los \emph{workers} de cómputo pesado están diseñados para ejecutarse como
procesos independientes en máquinas separadas que conectan con el RabbitMQ
central, reciben los trabajos, los procesan y escriben los resultados
directamente en la base de datos. Esto incluye tres modalidades ya
contempladas en el diseño actual:

\begin{enumerate}[leftmargin=1.5em]
  \item \textbf{Equipos propios del investigador} (GPU local disponible durante
    el trabajo diario).
  \item \textbf{Despliegue bajo SLURM en centros HPC externos} —por ejemplo,
    CESGA (Centro de Supercomputación de Galicia) vía sus convocatorias de
    acceso competitivo. Los \emph{workers} se empaquetan como imagen
    compatible (Singularity/Apptainer) y se lanzan como \emph{jobs} SLURM que,
    una vez activos, se conectan al RabbitMQ central. El número de
    \emph{workers} se escala pidiendo más \emph{jobs} SLURM, sin tocar el
    plano de control.
  \item \textbf{Otras infraestructuras \emph{cloud} o Kubernetes} (EOSC, BSC
    Life Sciences, nodos institucionales) bajo el mismo modelo.
\end{enumerate}

Esto significa que la contribución de una infraestructura \emph{host} —una
máquina modesta con alojamiento persistente y acceso externo— es precisamente
el recurso que los centros HPC no pueden ofrecer, y viceversa. Ambos roles son
complementarios y la arquitectura los coordina sin reconfiguración.

\subsection{Encaje con LifeWatch ERIC}

LifeWatch ERIC sostiene la infraestructura de investigación europea para
biodiversidad y ecosistemas. PROTEA cubre una necesidad concreta poco atendida
en ese ámbito: la anotación funcional de proteomas de organismos no-modelo,
donde las herramientas clásicas pierden potencia por la escasa representación
de estas especies en las bases de datos de referencia. Grupos de biodiversidad
que trabajan con proteomas recién ensamblados —metagenómica, genomas nuevos,
transcriptomas de especies no caracterizadas— obtienen con PROTEA una
anotación robusta sin necesidad de conocimiento técnico ni infraestructura
computacional propia.

El modelo de colaboración propuesto es escalonado y con riesgo mínimo en la
fase inicial: (1) fase piloto de 1–2 meses con despliegue y validación
extremo a extremo; (2) fase estable de 3–6 meses con apertura a grupos
externos y monitorización básica; (3) fase de integración de 6+ meses con
listado en el catálogo correspondiente y exploración de integración en los
entornos de trabajo científico del consorcio.

% =============================================================================
\section{Plan de publicación}
% =============================================================================

El trabajo prevé la publicación de un manuscrito de \emph{método} (predicción
de función vía re-ranker y evaluación CAFA~6) y otro de \emph{ingeniería}
(arquitectura de la plataforma), además de la eventual coautoría en el paper
retrospectivo del consorcio CAFA~6. Las revistas objetivo concretas, el orden
y el calendario se discutirán en un documento aparte.

% =============================================================================
\section{Cronograma orientativo}
% =============================================================================

\begin{itemize}[leftmargin=1.5em]
  \item \textbf{Corto plazo (0–3 meses)}: reconstrucción limpia y reproducible
    del método que obtuvo la posición CAFA~6 y validación del despliegue
    \texttt{docker compose} en entorno externo al equipo de desarrollo.
  \item \textbf{Medio plazo (3–12 meses)}: primer despliegue piloto en
    infraestructura externa (LifeWatch, si la vía se materializa, o
    alternativa comparable) y \emph{benchmarks} de reproducibilidad y
    \emph{throughput} de los pilares arquitectónicos.
  \item \textbf{Largo plazo (12–24 meses)}: consolidación del despliegue
    institucional e integración del trabajo en la memoria de tesis.
\end{itemize}

El cronograma es orientativo y se ajustará en función del cierre oficial del
\emph{ranking} CAFA~6 y de la disponibilidad real de infraestructura de
despliegue.

% =============================================================================
\section{Colaboradores y autoría}
% =============================================================================

El proyecto está co-dirigido por la Dra.~Ana~Rojas y el Dr.~David~Orellana,
ambos co-directores de la tesis, quienes proporcionan supervisión científica,
revisión crítica del trabajo y orientación sobre las vías de despliegue
institucional. La autoría final del paper se decidirá en función de las
contribuciones efectivas de cada parte al manuscrito, siguiendo las directrices
ICMJE.

Se agradece explícitamente el legado intelectual de la plataforma FANTASIA,
desarrollada previamente en el mismo ámbito, de la cual PROTEA hereda
conceptualmente parte de su enfoque pero sobre la que se ha realizado un
rediseño arquitectónico completo para resolver problemas de acoplamiento y
reproducibilidad heredados.

% =============================================================================
\section{Cierre}
% =============================================================================

PROTEA se encuentra en un momento en el que la parte técnica y metodológica
está sustancialmente madura y los resultados apuntan a una contribución
significativa en el campo de la anotación funcional de proteínas. Los
siguientes pasos —consolidación de los \emph{benchmarks}, despliegue en
infraestructura estable, redacción de los manuscritos y envío a las revistas
objetivo— son trabajo de los próximos 12–18 meses y constituyen el núcleo de
la fase final del doctorado.

Este documento se entrega como base de discusión para la revisión interna y
para facilitar la presentación del proyecto a colaboradores externos —en
particular, a la infraestructura LifeWatch ERIC como posible \emph{host} del
servicio en su fase de despliegue institucional.

\bigskip

\hrule

\vspace{1em}

\begin{thebibliography}{99}

\bibitem{cockburn2005hexagonal}
A.~Cockburn.
\newblock \emph{Hexagonal architecture}.
\newblock Artículo original disponible en alistair.cockburn.us, 2005.

\bibitem{martin2017clean}
R.~C.~Martin.
\newblock \emph{Clean Architecture: A Craftsman's Guide to Software Structure
  and Design}.
\newblock Prentice Hall, 2017.

\bibitem{ashburner2000go}
M.~Ashburner et al.
\newblock Gene Ontology: tool for the unification of biology.
\newblock \emph{Nature Genetics}, 25(1):25--29, 2000.

\bibitem{zhou2019cafa3}
N.~Zhou et al.
\newblock The CAFA challenge reports improved protein function prediction and
  new functional annotations for hundreds of genes through experimental
  screens.
\newblock \emph{Genome Biology}, 20(1):244, 2019.

\bibitem{lin2023esm2}
Z.~Lin et al.
\newblock Evolutionary-scale prediction of atomic-level protein structure with
  a language model.
\newblock \emph{Science}, 379(6637):1123--1130, 2023.

\bibitem{heinzinger2023prostt5}
M.~Heinzinger et al.
\newblock ProstT5: bilingual language model for protein sequence and
  structure.
\newblock \emph{bioRxiv}, 2023.

\bibitem{elnaggar2023ankh}
A.~Elnaggar et al.
\newblock Ankh: optimized protein language model unlocks general-purpose
  modelling.
\newblock \emph{bioRxiv}, 2023.

\bibitem{jones2014interproscan}
P.~Jones et al.
\newblock InterProScan 5: genome-scale protein function classification.
\newblock \emph{Bioinformatics}, 30(9):1236--1240, 2014.

\bibitem{cantalapiedra2021eggnog}
C.~P.~Cantalapiedra et al.
\newblock eggNOG-mapper v2: functional annotation, orthology assignments, and
  domain prediction at the metagenomic scale.
\newblock \emph{Molecular Biology and Evolution}, 38(12):5825--5829, 2021.

\bibitem{toronen2018pannzer2}
P.~Törönen et al.
\newblock PANNZER2: a rapid functional annotation web server.
\newblock \emph{Nucleic Acids Research}, 46(W1):W84--W88, 2018.

\bibitem{kulmanov2020deepgoplus}
M.~Kulmanov and R.~Hoehndorf.
\newblock DeepGOPlus: improved protein function prediction from sequence.
\newblock \emph{Bioinformatics}, 36(2):422--429, 2020.

\end{thebibliography}

\end{document}
